
\documentclass{article}
\usepackage{colortbl}
\usepackage{makecell}
\usepackage{multirow}
\usepackage{supertabular}

\begin{document}

\newcounter{utterance}

\centering \large Interaction Transcript for game `eqbench', experiment `eq\_bench', episode 83 with reflexion{-}llama{-}3.1{-}8b{-}instruct.
\vspace{24pt}

{ \footnotesize  \setcounter{utterance}{1}
\setlength{\tabcolsep}{0pt}
\begin{supertabular}{c@{$\;$}|p{.15\linewidth}@{}p{.15\linewidth}p{.15\linewidth}p{.15\linewidth}p{.15\linewidth}p{.15\linewidth}}
    \# & \multicolumn{2}{c}{Player} && \multicolumn{2}{c}{Game Master} \\
    \hline

    \theutterance \stepcounter{utterance}  
    & & & \multicolumn{4}{p{0.6\linewidth}}{
        \cellcolor[rgb]{0.9,0.9,0.9}{
            \makecell[{{p{\linewidth}}}]{
                \texttt{\tiny{[P1$\langle$GM]}}
                \texttt{You have attempted this same task before and did not succeed. These are the plans you wrote after those attempts.} \\
\texttt{Trial \#1: New plan: } \\
\\ 
\texttt{Upon re{-}examining the dialogue, I realize that I focused too much on the surface{-}level exchange between Charlie and Mari, without considering the underlying dynamics and power struggles at play. To better predict Mari's emotional response, I should have paid closer attention to the following:} \\
\\ 
\texttt{1. Mari's tone and language: Mari's responses are assertive and confident, with a hint of sarcasm. This suggests that Mari is not feeling defensive or guilty, but rather unapologetic and possibly even triumphant.} \\
\texttt{2. Charlie's behavior: Charlie's aggressive and condescending tone towards Mari implies that he is trying to assert his authority and control over the situation. This could be triggering Mari's unapologetic and triumphant emotions, as she is pushing back against Charlie's attempts to belittle her.} \\
\texttt{3. Mari's past: Mari mentions that she has been drinking herself into oblivion, which suggests that she may have a history of struggling with self{-}doubt and low self{-}esteem. However, in this dialogue, she is standing up to Charlie and refusing to back down, which indicates that she is feeling more confident and unapologetic.} \\
\\ 
\texttt{Given these observations, I would revise my scores as follows:} \\
\\ 
\texttt{Defensive: 2} \\
\texttt{Unapologetic: 9} \\
\texttt{Guilty: 0} \\
\texttt{Triumphant: 8} \\
\\ 
\texttt{Trial \#2: New plan: } \\
\\ 
\texttt{Upon re{-}examining the dialogue, I realize that I focused too much on the surface{-}level exchange between Charlie and Mari, without considering the underlying dynamics and power struggles at play. To better predict Mari's emotional response, I should have paid closer attention to the following:} \\
\\ 
\texttt{1. Mari's tone and language: Mari's responses are assertive and confident, with a hint of sarcasm. This suggests that Mari is not feeling defensive or guilty, but rather unapologetic and possibly even triumphant.} \\
\texttt{2. Charlie's behavior: Charlie's aggressive and condescending tone towards Mari implies that he is trying to assert his authority and control over the situation. This could be triggering Mari's unapologetic and triumphant emotions, as she is pushing back against Charlie's attempts to belittle her.} \\
\texttt{3. Mari's past: Mari mentions that she has been drinking herself into oblivion, which suggests that she may have a history of struggling with self{-}doubt and low self{-}esteem. However, in this dialogue, she is standing up to Charlie and refusing to back down, which indicates that she is feeling more confident and unapologetic.} \\
\\ 
\texttt{Given these observations, I would revise my scores as follows:} \\
\\ 
\texttt{Defensive: 2} \\
\texttt{Unapologetic: 9} \\
\texttt{Guilty: 0} \\
\texttt{Triumphant: 8} \\
\\ 
\texttt{To improve my prediction skills, I will pay closer attention to the underlying dynamics and power struggles in the dialogue, as well as the character's tone, language, and past experiences. I will also consider the context and setting of the dialogue, as well as any potential biases or assumptions that may be influencing my predictions.} \\
\\ 
\\ 
\\ 
\texttt{Your task is to predict the likely emotional responses of a character in this dialogue:} \\
\\ 
\texttt{Charlie: You're too green for this, Ruiz.} \\
\texttt{Mari: And you're too washed up, Redburn.} \\
\texttt{Charlie: I've been on this beat longer than you've been alive, kid.} \\
\texttt{Mari: And yet, here we are, stuck on the same case. Seems like experience isn't everything.} \\
\texttt{Charlie: Don't get smart with me, Mari.} \\
\texttt{Mari: Someone has to be the smart one here.} \\
\texttt{Charlie: You think I don't know what you're doing? Trying to prove yourself, make a name?} \\
\texttt{Mari: Better than drinking myself into oblivion.} \\
\texttt{Charlie: That's a low blow, Ruiz.} \\
\texttt{{[}End dialogue{]}} \\
\\ 
\texttt{At the end of this dialogue, Mari would feel...} \\
\texttt{Defensive} \\
\texttt{Unapologetic} \\
\texttt{Guilty} \\
\texttt{Triumphant} \\
\\ 
\texttt{Give each of these possible emotions a score from 0{-}10 for the relative intensity that they are likely to be feeling each.} \\
\\ 
\texttt{You must output in the following format, including headings (of course, you should give your own scores), with no additional commentary:} \\
\\ 
\texttt{Defensive: <score>} \\
\texttt{Unapologetic: <score>} \\
\texttt{Guilty: <score>} \\
\texttt{Triumphant: <score>} \\
\\ 
\\ 
\texttt{{[}End of answer{]}} \\
\\ 
\texttt{Remember: zero is a valid score, meaning they are likely not feeling that emotion. You must score at least one emotion > 0.} \\
\\ 
\texttt{Your answer:} \\
            }
        }
    }
     \\ \\

    \theutterance \stepcounter{utterance}  
    & \multicolumn{4}{p{0.6\linewidth}}{
        \cellcolor[rgb]{0.561,0.737,0.561}{
            \makecell[{{p{\linewidth}}}]{
                \texttt{\tiny{[P1$\rangle$GM]}}
                \texttt{Defensive: 2} \\
\texttt{Unapologetic: 9} \\
\texttt{Guilty: 0} \\
\texttt{Triumphant: 8} \\
            }
        }
    }
    & & \\ \\

    \theutterance \stepcounter{utterance}  
    & & & \multicolumn{2}{p{0.3\linewidth}}{
        \cellcolor[rgb]{0.9,0.9,0.9}{
            \makecell[{{p{\linewidth}}}]{
                \texttt{\tiny{[GM$|$GM]}}
                \texttt{Defensive: 2} \\
\texttt{Unapologetic: 9} \\
\texttt{Guilty: 0} \\
\texttt{Triumphant: 8} \\
            }
        }
    }
    & & \\ \\

    \theutterance \stepcounter{utterance}  
    & & & \multicolumn{2}{p{0.3\linewidth}}{
        \cellcolor[rgb]{0.9,0.9,0.9}{
            \makecell[{{p{\linewidth}}}]{
                \texttt{\tiny{[GM$|$GM]}}
                \texttt{Defensive: 4} \\
\texttt{Unapologetic: 8} \\
\texttt{Guilty: 0} \\
\texttt{Triumphant: 6} \\
            }
        }
    }
    & & \\ \\

    \theutterance \stepcounter{utterance}  
    & & & \multicolumn{2}{p{0.3\linewidth}}{
        \cellcolor[rgb]{0.9,0.9,0.9}{
            \makecell[{{p{\linewidth}}}]{
                \texttt{\tiny{[GM$|$GM]}}
                \texttt{game\_result = LOSE} \\
            }
        }
    }
    & & \\ \\

    \theutterance \stepcounter{utterance}  
    & & \multicolumn{4}{p{0.6\linewidth}}{
        \cellcolor[rgb]{0.95,0.95,0.95}{
            \makecell[{{p{\linewidth}}}]{
                \texttt{\tiny{[P1$\leftrightarrow$P]}}
                \texttt{Post{-}episode Reflexion (not sent to GM):} \\
\texttt{New plan: } \\
\\ 
\texttt{Upon re{-}examining the dialogue, I realize that I focused too much on the surface{-}level exchange between Charlie and Mari, without considering the underlying dynamics and power struggles at play. To better predict Mari's emotional response, I should have paid closer attention to the following:} \\
\\ 
\texttt{1. Mari's tone and language: Mari's responses are assertive and confident, with a hint of sarcasm. This suggests that Mari is not feeling defensive or guilty, but rather unapologetic and possibly even triumphant.} \\
\texttt{2. Charlie's behavior: Charlie's aggressive and condescending tone towards Mari implies that he is trying to assert his authority and control over the situation. This could be triggering Mari's unapologetic and triumphant emotions, as she is pushing back against Charlie's attempts to belittle her.} \\
\texttt{3. Mari's past: Mari mentions that she has been drinking herself into oblivion, which suggests that she may have a history of struggling with self{-}doubt and low self{-}esteem. However, in this dialogue, she is standing up to Charlie and refusing to back down, which indicates that she is feeling more confident and unapologetic.} \\
\\ 
\texttt{Given these observations, I would revise my scores as follows:} \\
\\ 
\texttt{Defensive: 2} \\
\texttt{Unapologetic: 9} \\
\texttt{Guilty: 0} \\
\texttt{Triumphant: 8} \\
\\ 
\texttt{To improve my prediction skills, I will pay closer attention to the underlying dynamics and power struggles in the dialogue, as well as the character's tone, language, and past experiences. I will also consider the context and setting of the dialogue, as well as any potential biases or assumptions that may be influencing my predictions.} \\
\\ 
\texttt{To further refine my approach, I will ask myself the following questions:} \\
\\ 
\texttt{* What are the power dynamics at play in the dialogue? Who has the upper hand, and how is Mari responding to Charlie's attempts to assert his authority?} \\
\texttt{* How does Mari's past experience and behavior influence her emotional response in the dialogue?} \\
\texttt{* What are the underlying themes and emotions that are driving the conversation between Charlie and Mari?} \\
\\ 
\texttt{By considering these factors and asking myself these questions, I hope to improve my ability to predict Mari's emotional response and provide more accurate answers in the future.} \\
            }
        }
    }
    & \\ \\

\end{supertabular}
}

\end{document}
